% ------------------------------------------------------------
% Page 94
% ------------------------------------------------------------

Malgré les grands feux de bois on ne pouvait faire remonter suffisamment le
thermomètre. Les mamans n’avaient d’autre alternative, pour réchauffer leurs enfants,
que de les mettre au lit, en plein jour, sous les édredons, quand on en possédait. Un
fois remis en état ces enfants pouvaient retourner dehors à leurs jeux, quitte à revenir
sous les couvertures lorsque le gel les atteignait à nouveau.

Véritablement se vivait là, une vie primitive, presque animale, faite de dépouillement,
de souffrance et de pauvreté. C’était semblable à ce que nous avions vu au Pompidou
dans certaines fermes.

Les pouvoirs publics consentaient un effort financier en faveur de ce hameau. Chaque
famille touchait annuellement une petite indemnité de maintien. Quand l’Inspecteur
des Eaux et Forêts passait pour compter le nombre exact de foyers intéressés par cette
mesure, on allumait des feux dans toutes les maisons abandonnées. Il s’agissait de
bénéficier au maximum de cette aide de l’Etat pour lutter contre la dépopulation des
montagnes. Le total à partager restait cependant bien modeste, et l’encouragement à
occuper les lieux bien relatif. On l’appréciait quand même !

Le culte et la fête apportaient à ces « braves gens » comme un rayon de soleil et
transformaient momentanément, en tous cas, leur horizon. Nous étions tous réjouis et
réconfortés de ce contact avec ces humbles paysans déshérités, mais qui ne se
plaignaient pas de leur sort. Au contraire, la satisfaction éclatait sur leurs visages, et
je ne crois pas que, en temps ordinaire, ils aient connu vraiment la tristesse ou
l’amertume.

Tous les plats qui constituaient le repas provenaient du dur labeur (y compris la pêche
ou la chasse) de ces hommes et de ces femmes sauf un : le plat de morue servi dans la
sauce blanche légèrement colorée. Je n’ai jamais rien dégusté de plus savoureux et
aujourd’hui je donnerai beaucoup pour en retrouver le goût. Voilà une recette
culinaire précieuse qui s’est probablement perdue... ! (mais non, vous trouverez celle de
Marianne Loupiac à la fin de ces pages)

Et puis cette vallée de la Brèze, relativement étroite en ses débuts, avec ses grands prés
séparant la route de la rivière, ses arbres majestueux, hêtres et érables, ses tournants
continuels suivis parfois d’une longue ligne droite, vallée de plus en plus abrupte à
mesure qu’on la pénètre, jusqu’à son dégagement par des lacets au bout desquels se
dévoile brusquement l’agglomération des Oubrets : quelle merveille de la nature ! Je
l’ai parcourue à nouveau récemment, confondu par ce spectacle.

\hfill Jean Gounelle
